\documentclass[../../../include/open-logic-section]{subfiles}

\begin{document}

\olfileid{sth}{replacement}{absinf}
\olsection{بدل رکھن دا اصول تے «مطلق لامتناہیت»}

ہُن اسیں \limofsize{} دی گل پچھے چھڈ کے بدل رکھن دے اصول دے (اندرلے) جواز دیاں کجھ ہور کوششاں ویکھاں گے، جیہڑیاں سیٹاں دے مرحلیاں وچ بنن دے خیال نوں سنجیدگی نال لیندیاں نیں۔

جدوں اسیں پہلی واری مرحلہ وار طریقے دا خاکہ دتا سی، اسیں کجھ اصول دتے سن جیہڑے سمجھاندے سن کہ ہر مرحلے اُتے کی ہُندا اے۔ ایہ \stageshier، \stagesord، تے \stagesacc سن۔ بعد وچ اسیں کجھ ہور اصول پائے جیہڑے مرحلیاں دی گنتی بارے دسدے سن۔ \stagessucc{} دسدا سی کہ سیٹ بنن دا کم کدی مکدا نہیں، تے \stagesinf{} دسدا سی کہ ایہ کم لامتناہی نمبر والے اک مرحلے توں وی لنگھدا اے۔

ایہ سوچنا واجب اے کہ ایہ پچھلے دو اصول کسے ہور کھلّھے اصول توں نکلدے نیں، جیویں کہ
\begin{enumerate}
	\item[] \stagesinex. مرحلے مطلق طور اُتے لامتناہی نیں۔ درجہ وار سلسلہ جِنّا اُچا ہو سکدا اے، اِنّا اُچا اے۔
\end{enumerate}
صاف اے کہ ایہ اک غیر رسمی اصول اے۔ پر بھاویں ایہ سیٹ دے اکٹھا کرن والے مرحلہ وار تصور توں جھٹ \emph{لازم نہیں آؤندا}، اوہدے نال \emph{میل} ضرور کھاندا لگدا اے۔ گھٹ توں گھٹ، \limofsize دے اُلٹ، ایہ خیال قائم رکھدا اے کہ سیٹ مرحلہ وار بندے نیں۔

ہُن آس ایہ اے کہ \stagesinex{} نوں ورت کے اسیں بدل رکھن دے اصول دا جواز دے سکدے آں۔ آؤ ویکھیے کہ ایہ کم کیویں ہو سکدا اے۔

\olref[ordinals][idea]{sec} وچ اسیں ویکھیا کہ اک ایہو جیہی حُسنِ ترتیب بنانی سوکھی اے جیہڑی (سہج سمجھ مطابق) $\omega+\omega$ نال ہم ساخت ہونی چاہیدی اے۔ دوجے لفظاں وچ، اسیں سوکھ نال اک مرحلہ وار طریقہ سوچ سکدے آں جیہدی ترتیب دی قسم (سہج سمجھ مطابق) $\omega+\omega$ اے۔ سو جے اسیں \stagesinex منّ لیا اے، تاں سانوں ضرور منّنا چاہیدا اے کہ درجہ وار سلسلے دا گھٹ توں گھٹ $\omega+\omega$ واں مرحلہ، یعنی $V_{\omega+\omega}$، اے۔ کیوں جے سلسلہ ایتھوں تک ضرور چل \emph{سکدا} اے۔

ایس خیال نوں عام کر کے ایویں آکھ سکدے آں۔ جیہڑی وی حُسنِ ترتیب ہووے، مرحلہ وار درجہ بندی بنن دا کم گھٹ توں گھٹ اوہدی حد تک چلنا چاہیدا اے۔ تے ایہدی ضمانت اسیں بس \olref[ordinals][ordtype]{thmOrdinalRepresentation} نوں اک \emph{مسلّمہ} منّ کے دے سکدے آں۔ ایہ سانوں دسے گا کہ ہر حُسنِ ترتیب کسے فون نوئمان ترتیبی عدد نال ہم ساخت اے۔ کیوں جے ہر فون نوئمان ترتیبی عدد اپنے درجے دے برابر ہووے گا، \olref[ordinals][ordtype]{thmOrdinalRepresentation} سانوں دسے گا کہ جدوں وی اسیں اپنے سیٹ نظریے وچ کسے حُسنِ ترتیب دا بیان دے سکیے، سیٹ بنن دا مرحلہ وار کم اوہ حُسنِ ترتیب توں اگے ودھے گا۔

ایہ خیال ضرور \stagesinex توں نکلدا اک نتیجہ لگدا اے۔ افسوس، جے ساڈا مقصد ایس خیال توں بدل رکھن دا اصول کڈھنا اے، تاں اک سادہ فنی اڑچن اے۔ بدل رکھن دا اصول \olref[ordinals][ordtype]{thmOrdinalRepresentation} توں ٹھیک ٹھیک ودھ طاقتور اے۔ (ایہ گل \citet[\S13.2]{Potter2004} نے دسی اے۔ اسیں ایہنوں \olref[sth][replacement][finiteaxiomatizability]{sec} وچ ثابت کراں گے۔)

نتیجہ ایہ اے کہ جے اسیں \stagesinex{} نوں ایویں سمجھنا اے کہ اوہتوں بدل رکھن دا اصول نکلے، تاں ایہ \emph{بس} ایہ نہیں آکھ سکدا کہ درجہ وار سلسلہ ہر حُسنِ ترتیب توں اگے ودھدا اے۔ ایہنوں ایہتوں وڈا دعویٰ کرنا پئے گا۔ ایس کم لئی \cite{Shoenfield:AST} نے ایس خیال دا اک بڑا سہج، ہور طاقتور روپ پیش کیتا۔ درجہ وار سلسلہ کسے وی سیٹ نال \emph{ہم حد} نہیں۔\footnote{گوڈل نے وی ایہو جیہا خیال دتا لگدا اے۔ \citet[p.~223]{Potter2004} ویکھو۔ گوڈل تے \citeauthor{Shoenfield:AST} بارے بحث لئی \citet[90--5]{Incurvati2020} ویکھو۔} ذرا ہور تفصیل نال، جے $\tau$ سیٹاں نوں درجہ وار سلسلے دے مرحلیاں ول بھیجن والا نقشہ اے، تاں کسے وی سیٹ $A$ دا $\tau$ تھلّے عکس سارے سلسلے نوں نہیں گھیر سکدا۔ خاکے دی صورت وچ ایویں آکھو۔
\begin{enumerate}
	\item[] \stagescofin. جے $A$ اک سیٹ اے تے $\tau(x)$ ہر $x \in A$ لئی اک مرحلہ اے، تاں اک مرحلہ اے جیہڑا ہر $\tau(x)$ توں پچھوں آؤندا اے، سارے $x \in A$ لئی۔
\end{enumerate}
صاف اے کہ $\ZF$ توں \stagescofin دا ڈھکواں رسمی روپ ثابت ہُندا اے۔ اُلٹے پاسے اسیں غیر رسمی طور اُتے دلیل دے سکدے آں کہ \stagescofin{} بدل رکھن دے اصول دا جواز دیندا اے۔\footnote{\stagescofin دے کسے رسمی روپ توں بدل رکھن دا اصول ثابت کرنا ہور اوکھا ہووے گا، کیوں جے اکلّا $\Z$ مرحلیاں دی تعریف کرن جوگا طاقتور نہیں۔ ایس لئی صاف نہیں کہ \stagescofin نوں رسمی صورت کیویں دیئیے۔ اک طریقہ ایہ اے کہ $\LT$ دی کسے توسیع وچ کم کریے، جیویں \olref[sth][replacement][extrinsic]{sec} وچ دَسیا اے۔} منّ لو $(\forall x \in A)\lexists![y][\phi(x,y)]$۔ فیر ہر $x \in A$ لئی، $\sigma(x)$ اوہ $y$ ہووے جیہدے لئی $\phi(x,y)$، تے $\tau(x)$ اوہ مرحلہ ہووے جتھے $\sigma(x)$ پہلی واری بندا اے۔ \stagescofin مطابق اک مرحلہ $V$ اے جیہدے لئی $(\forall x \in A)\tau(x)\in V$۔ ہُن کیوں جے ہر $\tau(x) \in V$ تے $\sigma(x) \subseteq \tau(x)$، اسیں وکھ کرن دے اصول نال $ \Setabs{y
\in V}{(\exists x \in A)\sigma(x) = y} = \Setabs{y}{(\exists x \in
A)\phi(x,y)}$ حاصل کر سکدے آں۔

\begin{prob}
	\stagescofin{} نوں $\ZF$ وچ رسمی صورت دیو۔
\end{prob}

سو \stagescofin{} بدل رکھن دے اصول دا جواز دیندا اے۔ تے ایہ گھٹ توں گھٹ منّن جوگی گل اے کہ \stagesinex{} توں \stagescofin دا جواز ملدا اے۔ کیوں جے منّ لو \stagescofin{} نئیں چلدا۔ فیر سلسلہ کسے سیٹ~$A$ نال ہم حد اے، یعنی اک نقشہ $\tau$ اے جیہدے لئی ہر مرحلے~$S$ واسطے کوئی $x \in A$ اے جس لئی $S \in \tau(x)$۔ ایس صورت وچ سلسلے دی منّی ہوئی «مطلق لامتناہیت» نوں فڑن دا طریقہ ساڈے کول اے۔ $\tau$ نوں $A$ اُتے لاؤن نال جیہڑی اقدار دا سیٹ ملدی اے، اوہ سارے سلسلے نوں \emph{گھیر لیندا اے}۔ تے ایہ گل سلسلے دے «مطلق طور اُتے لامتناہی» ہون دے خیال نوں کمزور کردی اے۔ عکسِ نقیض لئیے تاں \stagesinex{} توں \stagescofin لازم آؤندا اے، جیہڑا اگوں بدل رکھن دے اصول دا جواز دیندا اے۔

ایہ بدل رکھن دے اصول دا اندرلا جواز دین دی واقعی آس والی کوشش اے۔ پر آخر ایہ کوشش کامیاب ہُندی اے یا نہیں، ایہ فیصلہ اسیں تہاڈے اُتے چھڈاں گے۔

\end{document}
